\begin{proposition}[a nonzero normal Euler number]\label{push:src:C.3.11}
Glue the two fixed boundary spheres of \(Y=\mathrm{dP}_6/\rho\)
using their actual conductor identification, and call the resulting
closed oriented three-manifold \(N\). Its natural map to the
smooth cusp orbit five-manifold is an immersion. Its sole
self-intersection circle has two disjoint preimage circles, the
two pairs of nonfixed edge intervals of \ref{push:src:C.3.7}.
Along the fixed sphere \(F\subset N\), the oriented normal
two-plane bundle of this immersion has Euler number of absolute
value two. If its complex orientation is that of the Hopf invariant
\(a=xz\), where \(x,z\) have circle weights \(+1,-1\), the
Euler number is \(-2\).
\end{proposition}
\begin{proof}
Away from the conductor the normalization quotient is embedded.
At an ordinary nonfixed double point the circle acts freely along
the double curve, and the two smooth three-dimensional branches
intersect transversely along its one-dimensional orbit quotient.
At a fixed double point the two normal coordinates have weights
\(+1,-1\). The Hopf quotient coordinates are
\[
 a=xz\in\mathbb C,\qquad
 b=(|x|^2-|z|^2)/2\in\mathbb R,
\]
together with the coordinate along the fixed curve. The two
halves of \(Y\) glue as the two signs of \(b\), so their
union is smooth. At a triple point the full base equation is
\(q=xyz\), with \(y\) tangent to \(F\) and invariant.
In the quotient it becomes \(q=ay\). Hence the central orbit
fibre is the union of the two smooth three-planes
\(\{a=0\}\) and \(\{y=0\}\), intersecting transversely
along the \(b\)-axis. This verifies immersion and transversality
at both previously distinguished vertices. The endpoint
identifications of \ref{push:src:C.3.7} join the two nonfixed double arcs into
one circle; its preimages in \(N\) are disjoint circles.

The two complex weight lines in \(N_{F/X}\) are each
\(\mathcal O_{\mathbb P^1}(-1)\). To see this with the actual
normalization, each is the normal line of one fixed edge in its
\(\mathrm{dP}_6\) branch, and each edge has self-intersection
\(-1\). This remains true at the two vertices because those
branches and their edge normals are smooth there. The circle
characters distinguish the lines globally, so no choice of a
splitting of an extension is required.

The Hopf invariant \(a=xz\) is valued in the tensor product
of these two weight lines. It therefore gives the complex normal
line of the branch \(\{a=0\}\) along \(F\) as
\[
 \mathcal O(-1)\otimes\mathcal O(-1)=\mathcal O(-2).
\]
At a triple point \(q=ay\), the coordinate \(a\) still is
a normal coordinate of that branch. The zero of \(y\) does
not alter its normal bundle; using \(q\) instead would lose
precisely this distinction. The stated Euler number is consequently
\(-2\) for the \(a\)-complex orientation and \(+2\) if that
normal orientation is reversed.

In particular this immersion cannot simply be compressed along
\(F\) to an immersed three-sphere while retaining its normal
two-plane bundle: extension over a bounding three-disc would force
its Euler pairing on the boundary sphere to vanish. This identifies
a required framing correction; it does not compute that correction
or a Haefliger invariant for the annular twist.
\end{proof}
